\documentclass[11pt]{article}
\usepackage[T1]{fontenc}
\usepackage[scaled=0.92]{helvet}
\renewcommand{\familydefault}{\sfdefault}
\usepackage[a4paper,margin=2cm]{geometry}
\usepackage[british]{babel}
\usepackage{amsmath,amssymb}
\usepackage{graphicx}
\usepackage{enumitem}
\usepackage{booktabs}
\usepackage{titlesec}
\usepackage{xcolor}
\usepackage{tikz}
\usetikzlibrary{arrows.meta,positioning,fit,backgrounds}
\usepackage{pgfgantt}
\usepackage{caption}
\usepackage{csquotes}
\usepackage[numbers,sort&compress]{natbib}
\renewcommand{\bibsection}{\par\vspace{2pt}\noindent\textbf{References}\par}
\setlength{\bibsep}{0pt}
\usepackage[hidelinks]{hyperref}


\begin{document}

\begin{center}
{\large\bfseries PQ-Belenios: AI-Assisted Formal Verification for the Post-Quantum Migration of Verifiable Electronic Voting}\\[2pt]
{\small Mukesh Tiwari, Fran\c{c}ois Dupressoir, Nirav Ajmeri, Xiyue Zhang, Kerstin Eder (University of Bristol), Kristian Gj{\o}steen (NTNU)}
\end{center}

\section{Vision}
Belenios tallies complex voting methods such as the Schulze method with mixnets. 
All of its security rests on classical cryptography, so ballot secrecy is exposed to 
harvest-now-decrypt-later attacks: ciphertexts on today's bulletin board could be decrypted 
by a future quantum computer. 


{\small
\bibliographystyle{abbrv}
\bibliography{ref}
}

\end{document}
